\documentclass[conference]{IEEEtran}
\IEEEoverridecommandlockouts

\usepackage{cite}
\usepackage{amsmath,amssymb,amsfonts}
\usepackage{graphicx}
\usepackage{textcomp}
\usepackage{xcolor}
\usepackage{listings}
\usepackage{pgfplots}
\usepackage{dblfloatfix}

\lstset{
    basicstyle=\ttfamily\scriptsize,   % smaller font
    numbers=left,
    numberstyle=\tiny,
    frame=single,
    breaklines=true,
    columns=flexible
}

\begin{document}

\title{
Loss functions exploration for training a DCGAN on the CIFAR-10 dataset
}

\author{
\IEEEauthorblockN{PG60301 Rodrigo Domingues}
\IEEEauthorblockA{
MEI -- Computer Vision and Image processing \\
University of Minho \\
Email: pg60301@alunos.uminho.pt}
}

\maketitle

\begin{abstract}


\end{abstract}



\section{Introduction}

In this work, multiple loss functions are explored using a common DCGAN architecture in order to study the differences between several GAN training strategies. By keeping the network architecture and training setup fixed while modifying only the loss function, it is possible to directly compare their impact on training dynamics and image generation performance. Additionally, the CIFAR-10 dataset is divided into subsets of different sizes to investigate how the amount of available training data influences the effectiveness of each approach.

The evaluated variants are BCE, WGAN-GP, Hinge, LSGAN, RaGAN, R1, Feature Matching (FM), Perceptual Loss, MSGAN, and 0-GP.

The experiments were conducted using the CIFAR-10 dataset, which contains 50,000 RGB images of size $32 \times 32$ belonging to ten different classes. In addition to the full dataset, subsets containing 30,000 and 10,000 images were also considered in order to study the impact of dataset size on training performance.

Each configuration was trained for 100 epochs and periodically evaluated using the Fréchet Inception Distance (FID). The evolution of both generator and discriminator losses was also analyzed. The main objective is to compare the behavior of different loss functions and identify which approaches achieve the best performance under the considered experimental setup.



\section{Methodology}

\subsection{Dataset}

The CIFAR-10 dataset was selected for this work due to several practical and experimental considerations. First, it contains a diverse set of image classes with significant visual differences, making it suitable for evaluating how different GAN loss functions adapt to a variety of image distributions. The dataset consists of ten balanced classes: airplane, automobile, bird, cat, deer, dog, frog, horse, ship, and truck.

Another motivation for choosing CIFAR-10 is its relatively small size compared to more complex image datasets. The complete training set contains 50,000 RGB images with a resolution of $32 \times 32$ pixels, allowing multiple GAN variants to be trained within a reasonable computational budget while still providing meaningful comparisons between different training objectives.

The dataset is balanced, containing 5,000 training images per class. To investigate the influence of dataset size on GAN performance, two additional subsets containing 30,000 and 10,000 images were created. In order to avoid introducing class imbalance as a confounding factor, all subsets were constructed using a stratified sampling procedure, preserving the original class distribution of the full dataset.

\subsection{Network Architecture}

A common DCGAN architecture was used throughout all experiments. The generator receives a 100-dimensional latent vector sampled from a normal distribution and transforms it into a $32\times32$ RGB image through a series of transposed convolutional layers. Batch normalization and ReLU activations are applied after each intermediate layer, while a Tanh activation is used at the output layer.

The discriminator performs the opposite task, receiving either a real image from the dataset or an image generated by the generator. Through a sequence of convolutional layers, the image is progressively reduced to a single output value representing the discriminator's assessment. LeakyReLU activations with a negative slope of 0.2 are used throughout the network, while batch normalization is applied to all convolutional layers except the first.

The architecture uses 64 feature maps in the first layer of both the generator and discriminator, following the standard DCGAN design. To ensure a fair comparison between loss functions, the same generator and discriminator architectures were used for all experiments whenever possible. The only exception was the Feature Matching (FM) variant, which required a modified discriminator capable of returning intermediate feature representations in addition to the final output.

All convolutional layers were initialized using a normal distribution with mean 0 and standard deviation 0.02, while batch normalization layers were initialized according to the original DCGAN recommendations. Training was performed using the Adam optimizer with a batch size of 128, learning rates of $2\times10^{-4}$ for both networks, and $\beta_1=0.5$.

\subsection{Loss Functions}

The main objective of this work is to compare different GAN
training objectives while keeping the network architecture
unchanged. A total of ten loss function variants were evaluated.

\begin{itemize}

\item \textbf{BCE (Binary Cross Entropy)}:
The standard GAN objective proposed in the original GAN
framework. The discriminator performs binary classification
between real and generated samples.

\item \textbf{WGAN-GP}:
Replaces the binary classification objective with the Wasserstein
distance and introduces a gradient penalty to improve training
stability and reduce issues such as mode collapse.

\item \textbf{Hinge Loss}:
Uses a margin-based objective that only penalizes predictions
within a predefined margin.

\item \textbf{LSGAN}:
Substitutes binary cross-entropy with a mean squared error
objective. The smoother gradients can improve convergence
during training.

\item \textbf{RaGAN (Relativistic Average GAN)}:
Rather than estimating whether a sample is real or fake
independently, the discriminator learns whether real samples
appear more realistic than generated samples on average.

\item \textbf{R1 Regularization}:
Adds a gradient penalty on real samples to regularize the
discriminator and encourage smoother decision boundaries.

\item \textbf{Feature Matching (FM)}:
The generator is encouraged to match intermediate feature
statistics extracted by the discriminator instead of focusing
solely on fooling it.

\item \textbf{Perceptual Loss}:
Introduces feature representations extracted from a pretrained
VGG16 network. The generator is optimized to produce images
that are perceptually similar to real samples.

\item \textbf{MSGAN (Mode Seeking GAN)}:
Adds a mode-seeking term that encourages different latent
vectors to generate visually distinct outputs, aiming to improve
sample diversity.

\item \textbf{0-GP}:
Applies a zero-centered gradient penalty to regularize the
discriminator and promote smoother gradients during training.

\end{itemize}

Although some of these methods require additional components
or modifications to the standard GAN objective, all variants
were implemented within the same DCGAN framework. This
ensures that observed performance differences can be primarily
attributed to the loss function rather than architectural changes.
The only exception is MSGAN, which requires a slightly
modified training procedure. During each training iteration,
two generated samples are produced from different latent
vectors for every real batch, allowing the mode-seeking term
to measure output diversity and encourage the generator to
produce distinct samples.



\subsection{Evaluation Metric}

The primary evaluation metric used in this work was the
Fréchet Inception Distance (FID). This metric compares the
feature distributions of real and generated images using a
pretrained Inception network. Lower FID values indicate that
the generated images are closer to the real data distribution.

FID was computed every 10 epochs using 5,000 generated
images and a fixed reference set of 5,000 real CIFAR-10
images. This allowed the evolution of image quality to be
monitored throughout training while ensuring consistent
comparisons across different loss functions.

In addition to the final FID score, the evolution of generator
and discriminator losses was also analyzed to study training
stability and convergence behavior.





\end{document}
